\clearpage
\fancyhead[L]{\small COMET.}
\hypertarget{native-476-middle}{}
\begin{center}
\rule{0.55\linewidth}{0.4pt}\\[2pt]
Astronomische Nachrichten Band XXIII. Nr. 543. Seite 225. 226. 1845 August 7.\\[2pt]
\rule{0.55\linewidth}{0.4pt}
\end{center}

\begin{center}Meridianbeobachtung des von Herrn \textsc{Colla} entdeckten Cometen in Göttingen.\end{center}

Herr Geheimer Hofrath \textsc{Gauss} hat mir folgende am Göttinger Meridiankreise gemachte Beobachtung des Cometen in der unteren Culmination gefälligst mitgetheilt.

\begin{center}
\setlength{\tabcolsep}{5pt}
\begin{tabular}{l|r|r|r}
 & Mittlere Zeit & AR. & Nordl. Abw. \\
\hline
1845 Juni 12 & $13^{\rm u}\ 12^{\rm m}\ 59^{\rm s}9$ & $99^\circ\ 26'\ 5''4$ & $43^\circ\ 56'\ 53''6$ \\
\end{tabular}
\end{center}
\begin{flushright}\textsc{Schumacher}.\end{flushright}

\clearpage
\fancyhead[L]{\small NEPTUN.}
\hypertarget{native-476-planet}{}
\begin{center}
\rule{0.55\linewidth}{0.4pt}\\[2pt]
Astronomische Nachrichten Band XXV. Nr. 581. Beilage. Seite 82. 1846 October 22.\\[2pt]
\rule{0.55\linewidth}{0.4pt}
\end{center}

\begin{center}Erste Beobachtungen von \textsc{Le Verrier}'s Planeten.\end{center}

Herr Geheimer Hofrath \textsc{Gauss} hat mir folgende Meridianbeobachtungen gesandt.

\begin{center}
\setlength{\tabcolsep}{5pt}
\begin{tabular}{lr|r|r|r}
\multicolumn{2}{l|}{} & M. Z. in Göttingen & & \\
\hline
Sept. & 27 & $9^{\rm u}\ 27^{\rm m}\ 57^{\rm s}1$ & $328^\circ\ 14'\ 35''8$ & $-13^\circ\ 25'\ 54''0$ \\
Octob. & 6 & $8\ 51\ 56.5$ & $5\ 12.6$ & $-\phantom{13^\circ\ }29\ 6.0$ \\
 & 10 & $8\ 35\ 58.8$ & $1\ 40.7$ & $-\phantom{13^\circ\ }30\ 17.8$ \\
\end{tabular}
\end{center}
\begin{flushright}\textsc{Schumacher}.\end{flushright}

\clearpage
\hypertarget{native-476-letter}{}
\begin{center}
\rule{0.55\linewidth}{0.4pt}\\[2pt]
Astronomische Nachrichten. Band XXV. Nr. 582. Seite 95. 1846 November 5.\\[2pt]
\rule{0.55\linewidth}{0.4pt}
\end{center}
\begin{flushright}Göttingen 1846 October 22.\end{flushright}

Gestern Abend habe ich den \textsc{Le Verrier}'schen Planeten am Meridiankreise beobachtet, und ich verfehle nicht, Ihnen das Resultat mitzutheilen:

\begin{center}
1846 Octob. 21.\quad $7^{\rm u}\ 52^{\rm m}\ 13^{\rm s}7$\par
G. A. $=327^\circ\ 54'\ 7''8$\qquad Abw. $-13^\circ\ 32'\ 50''1$
\end{center}

Die Beobachtungen stimmten noch immer sämmtlich sehr gut mit einer Kreisbahn.
